\label{ch:gauss}

\section{Cofinalité entre cylindres HMT et cylindres de fractions continues}

Une section HMT irrationnelle détermine une famille emboîtée de cylindres compacts. Son évaluation archimédienne est un point \(x\in(0,1)\setminus\mathbb Q\). La fraction continue de \(x\) détermine des cylindres classiques

\[
I(a_1,\ldots,a_n).
\]

\begin{proposition}[Cofinalité]
Pour tout cylindre fini \(I(a_1,\ldots,a_n)\) qui contient l'évaluation \(x\), il existe une profondeur HMT \(m(n)\) telle que le cylindre HMT de niveau \(m\ge m(n)\) soit contenu dans \(I(a_1,\ldots,a_n)\).
\end{proposition}

\begin{proof}
\(n\) étant fixé, le cylindre de fraction continue \(I(a_1,\ldots,a_n)\) est un voisinage ouvert de l'irrationnel \(x\). Il existe alors \(\varepsilon_n>0\) avec \((x-\varepsilon_n,x+\varepsilon_n)\subset I(a_1,\ldots,a_n)\). Les cylindres HMT \(C_m\) contiennent \(x\), sont emboîtés et satisfont \(\operatorname{diam}C_m\to0\). Choisissons \(m(n)\) avec \(\operatorname{diam}C_{m(n)}<\varepsilon_n\). Comme \(x\in C_{m(n)}\), tout point de ce cylindre est à une distance de \(x\) inférieure à \(\varepsilon_n\), donc \(C_{m(n)}\subset I(a_1,\ldots,a_n)\). L'emboîtement donne l'inclusion pour tout niveau ultérieur.
\end{proof}

La démonstration utilise la compacité, l'emboîtement et l'unicité de
l'évaluation. Le résultat établit le transport des approximations, mais
n'implique pas la conjugaison du décalage de mémoire TPK avec l'application de
Gauss sur toutes les fibres profondes.

La généalogie complète de cette réalisation est
\[
\APP\to\TRIT\to\TPK
\to (C_m,\rho_{m+1,m})_{m\ge0}
\to x=\bigcap_m C_m
\to \bigl(I(a_1,\ldots,a_n)\bigr)_{n\ge1}.
\]
Les cylindres HMT conservent la route, la phase et la retenue ; l'application terminale projette
ces coordonnées sur les préfixes de fraction continue. La cofinalité
démontre que cette projection converge vers le même point. Un entrelaceur
dynamique exigerait en outre un relèvement \(\widetilde T_G\) avec des
carrés commutatifs sur les fibres de valuation, propriété qui ne se
déduit pas de la cofinalité.

\section{Mesure de Gauss et observables de Khinchin et de Lévy}

L'application de Gauss \cite{Khinchin,IosifescuKraaikamp}

\[
T_G(x)=\frac1x-\left\lfloor\frac1x\right\rfloor
\]

préserve

\begin{equation}
d\mu_G(x)=\frac{dx}{\log2(1+x)}.
\label{eq:gauss-measure}
\end{equation}

La constante de Khinchin est définie par

\begin{equation}
\log K=\int_0^1\log a_1(x)\,d\mu_G(x).
\label{eq:khinchin}
\end{equation}

La distribution du premier chiffre se calcule directement :
\begin{equation}
\mu_G(a_1=k)
=\frac1{\log2}\int_{1/(k+1)}^{1/k}\frac{dx}{1+x}
=\log_2\!\left(\frac{(k+1)^2}{k(k+2)}\right).
\label{eq:gauss-digit-law}
\end{equation}
Par le théorème ergodique,
\begin{equation}
K=\prod_{k\ge1}
k^{\,\log_2((k+1)^2/[k(k+2)])}
=2.685452001065306445309714835\ldots
\label{eq:khinchin-product}
\end{equation}
pour \(\mu_G\)-presque toute section. Le produit ne s'applique pas à chaque point :
la section \(\varphi^{-1}\), par exemple, a tous ses chiffres égaux
à un et constitue un critère de réfutation immédiat d'une lecture point par point. Une
certification de \cref{eq:khinchin-product} tronque la somme des logarithmes et
borne la queue au moyen de l'estimation
\[
\mu_G(a_1=k)=O(k^{-2}),
\qquad
\sum_{k>N}\mu_G(a_1=k)\log k
=O\!\left(\frac{\log N+1}{N}\right).
\]

La constante de Lévy est

\begin{equation}
L=-\int_0^1\log x\,d\mu_G(x)
=\frac{\pi^2}{12\log2}.
\label{eq:levy}
\end{equation}

En effet,

\[
\frac1{1+x}=\sum_{n\ge0}(-1)^nx^n,\qquad
\int_0^1(-\log x)x^n\,dx=\frac1{(n+1)^2},
\]

de sorte que l'intégrale est \(\eta(2)/\log2=\pi^2/(12\log2)\).

La constante de Khinchin détermine la moyenne logarithmique des quotients
partiels, tandis que la constante de Lévy détermine le taux exponentiel
de croissance des dénominateurs. Ce sont deux observables distinctes d'un
même système mesuré.

\section{Entropie métrique, exposant de Lyapunov et constante de Lochs}

Comme

\[
\log|T_G'(x)|=-2\log x,
\]

l'exposant de Lyapunov et l'entropie métrique sont

\begin{equation}
\chi_G=h_G=2L=\frac{\pi^2}{6\log2}.
\label{eq:gauss-entropy}
\end{equation}

Le rapport de Lochs entre les fractions continues et un développement en base \(b\) est

\begin{equation}
\Lambda_b=\frac{\log b}{h_G}
=\frac{6\log2\log b}{\pi^2}.
\label{eq:lochs}
\end{equation}

En particulier,

\begin{align*}
\Lambda_{10}&=0.9702701143920339257402560192\ldots,\\
\Lambda_{729}&=\frac{36\log2\log3}{\pi^2}
=2.777618966374305777705782976\ldots.
\end{align*}

La base \(729=3^6\) est pertinente pour l'espace ambiant \(\mathbb F_3^6\), mais \cref{eq:lochs} reste un théorème presque sûr de \(\mu_G\) ; et non une égalité point par point pour toute route HMT.

La valeur \(\Lambda_{729}\) représente le taux asymptotique typique de
quotients partiels obtenus par bloc complet de six trits. Elle n'identifie pas
l'espace ambiant de \(729\) mots aux autres objets de cardinal \(729\) du
corpus. L'égalité est un taux d'information entre partitions, et non une
bijection entre états.

\section{Opérateur de transfert de Gauss--Kuzmin--Wirsing}

L'opérateur de Gauss--Kuzmin--Wirsing agit par

\begin{equation}
(\mathcal L f)(x)=\sum_{n\ge1}\frac1{(x+n)^2}
f\!\left(\frac1{x+n}\right).
\label{eq:gkw-operator}
\end{equation}

Sa valeur propre dominante correspond à la densité invariante. La constante GKW est le module de la valeur propre sous-dominante sur le sous-espace de moyenne nulle. Un schéma HMT compatible utilise des partitions de \(9^m\) cylindres, des matrices de transfert avec intervalles et une borne de queue.

Pour promouvoir le calcul, il faut certifier : l'isolement de la valeur propre, l'invariance du sous-espace, l'erreur de troncature, la convergence des projecteurs spectraux et l'absence d'une autre valeur propre de même module. La valeur décimale connue ne peut pas être introduite comme centre de l'intervalle initial.

Une déformation conjointe contient Khinchin et Lévy dans un seul objet :
\begin{equation}
(\mathcal L_{s,t}f)(x)
=\sum_{n\ge1}n^t(n+x)^{-2s}
 f\!\left(\frac1{n+x}\right),
\qquad
P(s,t)=\log r(\mathcal L_{s,t}).
\label{eq:gauss-pressure-operator}
\end{equation}
Au point \((s,t)=(1,0)\), la perturbation spectrale donne
\begin{equation}
\partial_tP(1,0)=\log K,
\qquad
-\partial_sP(1,0)=\chi_G=2L.
\label{eq:pressure-derivatives}
\end{equation}
La hessienne de \(P\) détermine les variances et covariances asymptotiques de
\(\log a_1\) et \(-2\log x\). Par conséquent, les constantes de Khinchin et de
Lévy sont deux dérivées transversales de la même fonction de pression.

La reconnaissance numérique de la valeur propre sous-dominante est
\[
\lambda_2\simeq-0.3036630028987326586,\qquad
\kappa_{\rm GKW}=|\lambda_2|,\qquad
-\log\kappa_{\rm GKW}\simeq1.19183673556055.
\]
Dans ce volume, ces décimales constituent des valeurs de référence pour la
comparaison, et non un certificat spectral.
La preuve HMT proposée utilise Galerkin ou Ulam sur les cylindres cofinaux de
profondeur \(9^m\), une inégalité de Lasota--Yorke, l'arithmétique
d'intervalles pour isoler \(\lambda_2\) et une borne rigoureuse de la queue. Un
intervalle contenant deux valeurs propres de même module réfute la
promotion, sans pour autant affecter Khinchin, Lévy ou Lochs.

